% introduction.tex - Introduction to the Lambda^Cap white paper

\section{Introduction}
\label{sec:introduction}

Software systems increasingly execute code that is neither fully trusted nor
fully reviewed before deployment. %
Some of this code comes from plugins or user-defined scripts. %
Increasingly, it is generated on demand by large language models. %
In all of these settings, the security problem is the same: how do we ensure
that code receives exactly the authority it needs, and no more?

This is a \emph{least-privilege} problem: authority should be explicit, bounded,
and reviewable rather than ambient \cite{saltzer-schroeder}. %
Capability-based security gives the right general answer to least privilege and
confused-deputy problems \cite{hardy-confused}, %
but past capability models and capability-oriented languages such as E and Joe-E
leave several important problems unresolved \cite{e-lang,joe-e}. %

This paper presents a mechanized mathematical model of \emph{contextual
  capabilities}. %
The main purpose of the model is to provide a foundation for designing new
secure programming languages in which authority is explicit, statically tracked,
and controlled by program context. %

\subsection{Research Gap}

There are three major limitations in past capability models that affect the
usability of secure programming languages. %

First, a capability system should support \emph{static checking of capability
  usage}. %
It should be possible to read a function's type signature as a trustworthy upper
bound on the capabilities it may use. %
The object-capability model \cite{miller-robust,miller-paradigm} gives a strong
design discipline, %
but it does not by itself provide a compact type-level account of capability use
that can be checked compositionally. %

Second, a capability system should facilitate the \emph{propagation of
  fine-grained capabilities}. %
Capability-based programs do not carry a single coarse permission; they pass
user-scoped services, output channels, loggers, tool bindings, and other narrow
capabilities through higher-order code. %
Without language support, manual propagation through every interface quickly
becomes boilerplate-heavy. %
A useful model should make such propagation explicit and compositional, without
collapsing these capabilities into ambient authority. %

Third, a capability system should support \emph{simulation and testing} by
treating capabilities as abstract interfaces rather than hidden ambient
mechanisms. %
When capabilities remain explicit program values, the same program structure can
be exercised against restricted, attenuated, or mocked capabilities. %
Existing capability systems do not always make that form of substitution natural
at the level of the model. %

\subsection{Contextual Capabilities}

This paper proposes a mathematical model of \emph{contextual capabilities}. %
The basic idea is that capabilities are ordinary typed values that can be
explicitly provided or rebound for a context. %
Code deep in a call chain may then use them without every intermediate function
carrying them as ordinary parameters. %
Their use is tracked statically, but their availability is determined by
surrounding program context rather than by ambient authority. %

\subsection{Contribution}

We present \lambdaCC{}, a minimal calculus of contextual capabilities. %
Its contribution is to combine explicit capability values, static capability
tracking, and a higher-order operational account in one verified formalism. %
The resulting model supports reasoning about restricted execution,
capability-sensitive interfaces, authority attenuation, and the safe execution
of untrusted code. %
It is intended as a design model for new secure programming languages, rather
than as a complete implementation strategy by itself. %
